\chapter{Fase 7: la Torre en vivo}
\label{cap:d-f7}
\textbf{Fecha}: 2 de octubre de 2026. \textbf{Nota}: \codigo{20-torre-en-vivo.md}.

\textbf{Qué se buscaba}: que la campaña se ajuste cada semana mientras corre. Responde a dos incidentes: el de Big Data
(arquitectura de baja latencia para ajustar la campaña durante su ejecución) y el de Mercadotecnia (qué hacer si la
afluencia rebasa la capacidad).

\begin{opciones}[7.1 · Qué señales vigila · 2-oct-2026]
\textbf{Elegidas, las cuatro}: ocupación del norte (DataTur semanal), tormentas cerca (HURDAT2, con la regla de la Fase 2),
clima raro (Isolation Forest) y llegadas mensuales del sur contra lo esperado (el rango del 90\,\% del Pronóstico).

\textbf{Detalle}: llegadas \emph{por debajo} del rango no pausan, porque eso es espacio, lo que la campaña busca.
\end{opciones}

\begin{opciones}[7.2 · Qué hacer en el sur con una mala semana · 2-oct-2026]
\begin{center}
\small
\begin{tabularx}{\linewidth}{>{\raggedright\arraybackslash}p{0.3\linewidth}Y}
\encabezado \h{Opción} & \h{Por qué sí o por qué no} \\
\textbf{Pausar esa semana y guardar el dinero (elegida)} & El dinero pasa a la siguiente semana permitida del mismo lugar; no
se pierde \\
\filagris Bajar a la mitad & Se seguiría anunciando con mal clima \\
Solo avisar & No cuida nada \\
\bottomrule
\end{tabularx}
\normalsize
\end{center}
\end{opciones}

\begin{opciones}[7.3 · Qué hacer si el norte se llena · 2-oct-2026]
\textbf{Elegida}: \textbf{encender ``¿Ibas al norte?''} si Cancún o la Riviera Maya llegan a 85.92\,\% (el corte de
``saturado'' del Radar), hacia el lugar del sur encendido con más espacio ese mes. \textbf{Descartado}: solo marcarlo.
\end{opciones}

\begin{opciones}[7.4 · Cómo mostrarlo · 2-oct-2026]
\textbf{Elegida}: \textbf{reproducción grabada} en la página, porque GitHub Pages es estático y no puede transmitir.
\textbf{Descartado en ese momento}: además, transmisión en vivo con un servidor local. (Se agregó después, en la Fase 9.)

\textbf{Cómo funciona}: cada semana real se escribe como un archivo, y \textbf{Spark Structured Streaming} los lee de uno en
uno (\codigo{maxFilesPerTrigger = 1}), en orden, como si llegaran. En cada lote, el motor de reglas decide. Cambiar la carpeta
por una fuente real no cambia el motor.
\end{opciones}

\begin{opciones}[7.5 · Qué es ``clima raro'' · 2-oct-2026]
\textbf{Lo que pasó}: el corte automático de Isolation Forest marcaba como raras el \textbf{47\,\%} de las semanas de Chetumal (e
incluso 13–26\,\% dentro de su propio entrenamiento). Con 156 semanas no distinguía, y desde 2022 llueve más y hace más calor
que en 2019–2021. Pausaría media campaña.

\begin{center}
\small
\begin{tabularx}{\linewidth}{>{\raggedright\arraybackslash}p{0.3\linewidth}Y}
\encabezado \h{Opción} & \h{Por qué sí o por qué no} \\
\textbf{Más raro que el 95\,\% de las semanas que el bosque ya conoce (elegida)} & Con un bosque por año que aprende de todos
los años anteriores: 21 semanas raras en Chetumal (8.8\,\%) y 17 en Kohunlich (7.1\,\%) \\
\filagris Más raro que el 99\,\% & Casi nunca pausa \\
Quitar la señal & Se pierde la vigilancia del clima \\
\bottomrule
\end{tabularx}
\normalsize
\end{center}
\end{opciones}

\section*{Resultado: 239 semanas reales}
\begin{itemize}
\item \textbf{239 lotes de Spark, en orden}, del 3 de enero de 2022 al 27 de julio de 2026.
\item \textbf{48 pausas} de lugar-semana: 44 con clima raro, 9 con tormenta cerca (Lisa 2022, Nadine y Sara 2024) y 5 porque
  llegó más gente de la esperada a Chetumal en mayo de 2025 (algunas con más de un motivo).
\item \textbf{No se perdió dinero}: los \$876,901 planeados se gastaron todos.
\item \textbf{``¿Ibas al norte?'' se encendió 6 semanas}, todas hacia la Bahía. El norte casi nunca se satura en una semana
  suelta (Cancún 1, Riviera Maya 9): su señal fuerte es la temporada.
\item \textbf{Validación independiente}: el bosque marcó las tres semanas de tormenta sin saber que hubo tormenta.
\end{itemize}

\textbf{Limitaciones declaradas}: el sur no tiene ningún dato semanal oficial; las tormentas de 2026 todavía no se publican
(``sin dato'', no ``sin tormenta''); se supone que el dato del mes se conoce al empezar el siguiente (en realidad tarda 1–2
meses); el plan de dinero se aplicó a años pasados.
